\documentclass[a4paper,12pt]{article}
\usepackage[utf8]{inputenc}
\usepackage[T1]{fontenc}
\usepackage[spanish]{babel}
\usepackage{geometry}
\geometry{margin=2.4cm}

\title{Extracto formativo de contexto\\Ley 39/2015, de Procedimiento Administrativo Común}
\author{Material de apoyo para la sesion 3}
\date{}

\begin{document}

\maketitle

\section*{Uso del documento}

Este documento se prepara como contexto para una practica con Codex.

No sustituye a la consulta de la norma oficial ni a la revision juridica de una persona competente. Para cualquier uso real debe consultarse el texto oficial y actualizado en el Boletin Oficial del Estado.

Fuente oficial de referencia:

\begin{itemize}
    \item Ley 39/2015, de 1 de octubre, del Procedimiento Administrativo Comun de las Administraciones Publicas.
    \item BOE-A-2015-10565.
    \item Texto consolidado disponible en: https://www.boe.es/buscar/act.php?id=BOE-A-2015-10565
    \item En la consulta realizada para preparar este material, el BOE indicaba ultima actualizacion publicada el 06/11/2024.
\end{itemize}

\section*{Regla de trabajo para Codex}

Cuando se use este PDF como contexto, Codex debe:

\begin{itemize}
    \item usar solo la informacion contenida en este documento;
    \item no inventar articulos, plazos, jurisprudencia ni organos;
    \item diferenciar entre resumen orientativo y cita normativa;
    \item marcar como ``pendiente de verificacion juridica'' cualquier duda;
    \item recomendar la consulta del BOE para uso real.
\end{itemize}

\section*{Materias seleccionadas}

Este extracto se centra en cuestiones utiles para documentos administrativos sencillos:

\begin{itemize}
    \item derechos de las personas interesadas;
    \item obligacion de resolver;
    \item motivacion de actos administrativos;
    \item notificacion;
    \item subsanacion de solicitudes;
    \item tramite de audiencia;
    \item contenido de la resolucion;
    \item recursos administrativos.
\end{itemize}

\section*{Articulo 13. Derechos de las personas}

La Ley 39/2015 reconoce derechos de las personas en sus relaciones con las Administraciones Publicas. Entre ellos, se incluye el derecho a comunicarse por medios electronicos, a ser asistidas en el uso de medios electronicos, a utilizar las lenguas oficiales en el territorio correspondiente, a acceder a informacion publica, archivos y registros en los terminos previstos en la normativa, y a exigir responsabilidades cuando proceda.

Para una practica de redaccion administrativa, esto implica que un informe no debe tratar a la ciudadania como mera destinataria pasiva, sino como titular de derechos.

\section*{Articulo 21. Obligacion de resolver}

La Administracion esta obligada a dictar resolucion expresa y a notificarla en todos los procedimientos, cualquiera que sea su forma de iniciacion.

En un informe interno, si se habla de expedientes, solicitudes o reclamaciones, debe tenerse presente la necesidad de identificar el estado de tramitacion y evitar expresiones que sugieran cierre del asunto si no consta resolucion.

\section*{Articulo 35. Motivacion}

Determinados actos administrativos deben ser motivados. Entre otros, aquellos que limiten derechos subjetivos o intereses legitimos, los que resuelvan recursos administrativos, los que se separen del criterio seguido en actuaciones precedentes y los que se dicten en ejercicio de potestades discrecionales.

En una propuesta o informe, la motivacion exige explicar razones, hechos relevantes y fundamento de la decision. No basta con una conclusion generica.

\section*{Articulos 40 a 44. Notificacion}

La notificacion es una garantia esencial para que la persona interesada conozca el contenido del acto administrativo. La notificacion debe practicarse conforme a las reglas legales aplicables y debe permitir conocer el texto del acto y los recursos que procedan, cuando corresponda.

En documentos de trabajo, no debe afirmarse que una persona esta correctamente notificada si no consta la practica de la notificacion conforme al procedimiento aplicable.

\section*{Articulo 68. Subsanacion y mejora de la solicitud}

Cuando una solicitud no reuna los requisitos exigidos, se debe requerir a la persona interesada para que subsane la falta o acompañe los documentos preceptivos, en los terminos previstos por la ley.

En una comunicacion administrativa, conviene indicar de forma clara que debe corregirse, que documento falta y que consecuencia puede tener no atender el requerimiento.

\section*{Articulo 82. Tramite de audiencia}

El tramite de audiencia permite a las personas interesadas conocer el procedimiento antes de la propuesta de resolucion y formular alegaciones o aportar documentos.

En un informe, si se recomienda resolver un expediente, debe comprobarse si procede o no dar audiencia y si ese tramite ya se ha realizado.

\section*{Articulo 88. Contenido de la resolucion}

La resolucion que ponga fin al procedimiento debe decidir todas las cuestiones planteadas por las personas interesadas y aquellas otras derivadas del procedimiento.

Para la redaccion de documentos, esto exige evitar respuestas parciales cuando la solicitud o el expediente plantean varias cuestiones.

\section*{Articulos 112 a 124. Recursos administrativos}

La Ley 39/2015 regula los recursos administrativos, entre ellos el recurso de alzada y el recurso potestativo de reposicion.

En un borrador de resolucion o comunicacion, no deben inventarse recursos, plazos ni organos competentes. Si el documento de contexto no contiene el detalle suficiente, debe marcarse como pendiente de verificacion juridica.

\section*{Ejemplo de encargo para Codex}

Usa este PDF como unica fuente normativa.

Necesito revisar un borrador de informe municipal sobre una solicitud ciudadana incompleta.

Indica:
\begin{itemize}
    \item que cautelas procedimentales deben comprobarse;
    \item que informacion no puede darse por supuesta;
    \item que partes del texto requieren verificacion juridica;
    \item que articulos de este PDF justifican cada recomendacion.
\end{itemize}

No inventes jurisprudencia ni normativa no incluida en el PDF.

\end{document}
